% Nouveau EGA V (cote 25): archive cote 25, site page 42 (PDF page 43).
% typescript. Draft compared: none.
% AI-transcribed by Claude Opus 5.5 (claude-opus-5-5, adaptive thinking, effort high), Message Batches API, batch ega25b (msgbatch_01A7n5YuXubUPNK1zsepi8RG), 2026-11-13; original AI transcription, subsequently checked against both supplied scan parts; see PAGES_CHECKED.tsv and CHECK_DECISIONS.tsv.
\scanpage{25/42}
% NOTE: The raised supsi belongs to the cancelled Y as a superscript; the cancellation crosses the base expression. All diagram arrows are penned. A separate dense blot beside the upper0 has no secure character value.
\mspage{\hand{38}}

$\underline{Y}^{\mathrm{sing}} = P(\underline{\hand{\nu}}_{X/P}) \hand{\longrightarrow} \overset{\hand{\vee}}{P}$ . Alors H'est ``orthogonal'' à x, et en parti-%
culier si dim H' = r-1 (ce qui équivaut par b) au fait que \xout{yx/x$\underline{Y}^{\unsure{supsi}}$xxxxx}
y est point \inserted[typed]{géom.} singulier \uline{ordinaire} sur $Y_{\hand{\xi}}$ ) x est uniquement déterminé
en fonction de H'.

Dém . Nous avons déjà prouvé a), et b) en est une conséquence
grâce aux développements du par.16. Reste à prouver c), qui est
encore conséquence d'une relation de compatibilité essentiellement
triviale. Je me borne ici à copier, dans le cas spécial qui nous
occupe (mais avec S quelconqque), le diagramme (D') du par.16, qui
prend la forme suivante (où nous écrivons $\underline{Y}^1$ pour $\underline{Y}^{\mathrm{sing}}$):

\diagram{
\begin{tikzcd}[column sep=small]
 & & 0 & & \\
 & & \hand{\underline{\Omega}}^1_{X/S\,\underline{Y}^1} \arrow[u] & & \\
0 \arrow[r] & \underline{P}^1(\underline{O}(1,1))\otimes\underline{O}_{\underline{Y}^1} \arrow[r] \arrow[ur, "\hand{\unsure{\nu}.\mu}"] & (\hand{\underline{\Omega}}^1_{X/S})_{\underline{Y}^1} + (\hand{\underline{\Omega}}^1_{\overset{\hand{\vee}}{P}/S})_{\underline{Y}^1} \arrow[r] \arrow[u] & \hand{\underline{\Omega}}^1_{\underline{Y}^1/S} \arrow[r] & 0 \\
 & \underline{O}_{\underline{Y}^1}(-1,-1) \arrow[u] \arrow[r, "\hand{u}"] & \hand{\underline{\Omega}}^1_{\overset{\hand{\vee}}{P}/S\,\underline{Y}^1} \arrow[u] \arrow[ur, "\hand{\nu}"'] & & \\
\hand{\underline{\omega}}^{-2}_{X/Y}\otimes\underline{O}_{\underline{Y}^1}(-2,-2) \arrow[ur] & & \hand{0} \arrow[u] & &
\end{tikzcd}
}

La relation d'orthogonalité c) ne fait qu'exprimer que le composé
$\hand{\nu}u$ dans ce diagramme est nul; je me dispense de développer plus
avant ce deuxième exercice de symtaxe. Je suppose que les gens avaient
l'habitude de prendre des coordonnées, et ils y voyaient peut-être plus
clair ....

\uline{Corollaire} 8.\over{\illegible{1}}{8}. Supposons S=Spec(k), \xout{\illegible{1} sorte que $\underline{Y}^{\mathrm{sing}}$} \inserted{$f : X \to P$ une immersion} et X
irréductible \inserted{(tjrs lisse sur k)}. \xout{\unsure{Alors}} Alors $\underline{Y}^{\mathrm{sing}}$ est lisse et irréductible; soit T

\margin{\hand{42}}
